% The "blind" option will make anonymous all author, affiliation, correspondence and funding information.
% Use "num-refs" option for numerical citation and references style.
% Use "alpha-refs" option for author-year citation and references style.
\documentclass[alpha-refs]{wiley-article}

\usepackage{color}
\usepackage{ulem}

% Uncomment to display with annotation; comment out otherwise
\newcommand{\add}[1]{\textcolor{blue}{#1}}
\newcommand{\delete}[1]{\textcolor{red}{\sout{#1}}}
\newcommand{\edit}[2]{\textcolor{red}{\sout{#1}} \textcolor{blue}{#2}}
\newcommand{\mnote}[1]{\marginpar{\textcolor{green}{\textbf{#1}}}}

 % Uncomment to display without annotation; comment out otherwise
%\newcommand{\add}[1]{#1}
%\newcommand{\delete}[1]{}
%\newcommand{\edit}[2]{#2}
%\newcommand{\mnote}[1]{}

\newcommand{\todo}[1]{\textcolor{red}{(#1)}}

% Update article type if known
\papertype{Original Article}

\title{Processes in hailstorms in Australia's tropics}

% Include full author names and degrees, when required by the journal.
\author[1,2,3]{Timothy H. Raupach}
\author[4]{Sarah D. Bang}

% Include full affiliation details for all authors
\affil[1]{UNSW Sydney Institute for Climate Risk and Reponse, UNSW Sydney, Sydney, New South Wales, Australia}
\affil[2]{UNSW Sydney Climate Change Research Centre, UNSW Sydney, Sydney, New South Wales, Australia}
\affil[3]{ARC Centre of Excellence for 21st Century Weather, Sydney, New South Wales, Australia}
\affil[4]{NASA Marshall Space Flight Center, Huntsville, Alabama, United States of America}

\corraddress{T. H. Raupach, UNSW Institute for Climate Risk and Response, UNSW Sydney, Sydney, New South Wales, 2052, Australia}
\corremail{t.raupach@unsw.edu.au}

\fundinginfo{THR's position is supported by QBE Insurance. This work used compute time provided by the National Computational Infrastructure National Computational Merit Allocation Scheme 2025 (grant ID NCMAS-2025-10).}

% Include the name of the author that should appear in the running header
\runningauthor{Raupach et al.}

\begin{document}

\maketitle

\begin{abstract}
    \todo{Abstract}

% Please include a maximum of seven keywords
\keywords{hail, severe weather, hailstorms, tropical processes}
\end{abstract}

\section{Introduction}



\section{Data and Methods}

\subsection{Hailstorm detections}

Hail signatures were those detected using the algorithm of \cite{Bang_JAMC_2019}.

\subsection{Downscaled simulations}

We used the Advanced Research Weather and Forecasting model
\citep[AR-WRF,][]{Skamarock_2021} v4.6.0 to simulate the hours around each
satellite hailstorm detection. For each simulation, we ran WRF for a period of
15 hours of simulation time, comprising 12 hours of model spin-up which was not
considered in the analyses, the hour before the hail detection, the hour of the
detection, and the hour afterwards. Boundary conditions were provided by
European Centre for Medium-Range Weather Forecasts Reanalysis 5
\citep[ERA5,][]{Hersbach_QJRMS_2020} data. The models were run with three nested
domains of 120 $\times$ 120 mass points each, centred on the satellite hail
detection; with outer to inner domain grid spacings of 9 km, 3 km, and 1 km. 79
vertical mass points were used with a model-top pressure of 50 hPa. The model
timestep was 18 seconds and output was written at 5 minute intervals. The domain
with 9 km grid spacing had the New Tiedtke \citep{Zhang_JC_2017} cumulus scheme
enabled, while no cumulus parameterisation was used in the finer-resolution
domains. Across all domains, we used the RRTMG radiation scheme
\citep{Iacono_JGRA_2008}, the revised MM5 Monin-Obukhov surface layer scheme
\citep{Jimenez_MWR_2012}, the YSU planetary boundary layer scheme
\citep{Hong_MWR_2006}, and the Noah-MP land-surface scheme
\citep{Niu_JGRA_2011}. WRF-HAILCAST \citep{Adams-Selin_WF_2019} was enabled on
the 3 km and 1 km domains to estimate surface hailstone sizes. For each
hailstorm detection we ran four model instances only differing in the choice of
microphysics scheme. The four schemes considered were the NSSL 2-moment with
cloud condensation nucleii prediction \citep[][]{Mansell_JAS_2010}, Thompson
2-moment \citep{Thompson_MWR_2008}, P3 3-moment \citep[][]{Milbrandt_JAS_2021},
and MY2 \citep{Milbrandt_JAS_2005a, Milbrandt_JAS_2005b} schemes.

% \begin{table}[bt]
% \caption{}
% \begin{threeparttable}
% \begin{tabular}{lccrr}
% \headrow
% \thead{Variables} & \thead{JKL ($\boldsymbol{n=30}$)} & \thead{Control ($\boldsymbol{n=40}$)} & \thead{MN} & \thead{$\boldsymbol t$ (68)}\\
% Age at testing & 38 & 58 & 504.48 & 58 ms\\
% Age at testing & 38 & 58 & 504.48 & 58 ms\\
% Age at testing & 38 & 58 & 504.48 & 58 ms\\
% Age at testing & 38 & 58 & 504.48 & 58 ms\\
% \hiderowcolors
% stop alternating row colors from here onwards\\
% Age at testing & 38 & 58 & 504.48 & 58 ms\\
% Age at testing & 38 & 58 & 504.48 & 58 ms\\
% \hline  % Please only put a hline at the end of the table
% \end{tabular}
% \begin{tablenotes}
% \item WRF, Weather Research and Forecasting.
% \end{tablenotes}
% \end{threeparttable}
% \end{table}

\section{Results}

\subsection{Skew-T analyses}

\subsection{Hydrometeor profiles}

\subsection{}

\section{Discussion and Conclusion}

% \begin{figure}[bt]
% \centering
% %\includegraphics[width=6cm]{example-image-rectangle}
% \caption{Although we encourage authors to send us the highest-quality figures possible, for peer-review purposes we are can accept a wide variety of formats, sizes, and resolutions. Legends should be concise but comprehensive – the figure and its legend must be understandable without reference to the text. Include definitions of any symbols used and define/explain all abbreviations and units of measurement.}
% \end{figure}


\section*{Acknowledgements}
This research/project was undertaken with the assistance of resources and
services from the National Computational Infrastructure (NCI), which is
supported by the Australian Government. Since March 2024, THR's position at UNSW
has been financially supported by QBE Insurance.

%\section*{Conflict of interest} You may be asked to provide a conflict of
% interest statement during the submission process. Please check the journal's
% author guidelines for details on what to include in this section. Please
% ensure you liaise with all co-authors to confirm agreement with the final
% statement.

\printendnotes
\bibliography{library}

\end{document}
